\documentclass[11pt]{article}
\usepackage{fontspec,amsmath,amssymb,amsthm,hyperref}
\usepackage[margin=1in]{geometry}
\setlength{\emergencystretch}{3em}
\title{Full-potential PIPULA and its smoothed-center theorem}
\author{LangevinLean mathematical exposition}
\date{}
\begin{document}\maketitle
\noindent Source module: \texttt{ProximalLangevin/PIPULA.lean}. Compile this proof document with XeLaTeX or LuaLaTeX. The accompanying HTML pairs the mathematical argument with the exact Lean statements and records the scope of the formalization.


Assume the full nonsmooth potential \(U\) is \(m\)-strongly convex and bounded below, with \(m,\gamma>0\). Global Lipschitzness of \(U\) is not required. The theorem concerns the maximizer of the marginal of \(U^\gamma\).


\section{The actual full-envelope modulus makes the step admissible}


Set \(m_\gamma=m/(1+m\gamma)>0\). Since \(\gamma m_\gamma=m\gamma/(1+m\gamma)\le1\), we have \(\gamma(m_\gamma+\gamma^{-1})\le2\), equivalently \(\gamma\le2/(m_\gamma+\gamma^{-1})\). All multiplications and divisions have positive denominators.




\textbf{Correspondence with the displayed declarations.}


\begin{itemize}

\item \texttt{fullSmoothingModulus}: Strong convexity after smoothing the entire potential.
\item \texttt{pipula\_admissible\_step}: With smoothing equal to the time step, the envelope always meets the proved Gaussian ULA step restriction.
\end{itemize}


\noindent Lean declarations: \texttt{fullSmoothingModulus}, \texttt{fullSmoothingModulus\_pos}, \texttt{pipula\_admissible\_step}.

\section{The proximal/noise law is the actual IPLA law}


The law is the pushforward of the exact PIPULA recursion. Set \(g_1=0\) and \(\eta=\gamma\) in the pointwise MYIPLA identity. Actual proximal optimality gives \(P_\gamma=I-\gamma\nabla U^\gamma\), so the full chain equals IPLA for \(U^\gamma\), and hence their pushforwards agree.




\textbf{Correspondence with the displayed declarations.}


\begin{itemize}

\item \texttt{pipulaLaws}: The probability law defined by the paper's pairwise PIPULA recursion.
\item \texttt{pipulaLaws\_eq\_iplaLaws}: Exact equality between the displayed PIPULA law and IPLA for the genuine full Moreau envelope.
\end{itemize}


\noindent Lean declarations: \texttt{pipulaLaws}, \texttt{pipulaLaws\_eq\_iplaLaws}.

\section{Sharp concentration and Gaussian discretization}


The algebraic reduction \(g_1=0\), \(\lambda=\gamma\) is proved for the actual update and its full recursion. Together with the envelope theory, it supplies genuine proximity, smoothness, and strong-convexity facts for a full-potential scheme.



It does not reproduce v3 Corollary A.6's original-center finite-time bound. In assumption B1, the displayed proximal inner-product deficit \(\mu\) gives the envelope-gradient constant \(\mu/\gamma\), since the residual is divided by \(\gamma\). The code records that exact implication. A scale-uniform PIPULA discretization theorem and an original-center estimate with the paper's rate remain unproved here; substituting a \(1/\gamma\) smoothness bound into a generic ULA estimate alone would not establish the claimed vanishing error.



A separate proved finite-time result is supplied for the actual PIPULA chain around the maximizer of its \emph{smoothed} marginal. Assume the original finite-valued potential is \(m\)-strongly convex and bounded below; global Lipschitzness of the strongly convex potential is not imposed. For \(\gamma>0\), its full envelope has modulus \(m_\gamma=m/(1+m\gamma)\) and gradient Lipschitz constant \(1/\gamma\). The actual proximal/noise recursion equals IPLA for that envelope, and the admissible-step inequality holds automatically. The resulting sharp concentration and Gaussian-chain bound are fully proved. Its conservative discretization term is \((33/20)(1/(\gamma m_\gamma))\sqrt{d_x+d_\theta/N}\sqrt\gamma\), whose coefficient diverges as \(\gamma\downarrow0\). This is not a reproduction of the paper's original-center \(O(\sqrt\gamma)\) convergence rate.



Writing \(m_\gamma=m/(1+m\gamma)\), \(\Pi_{N,\gamma}\) for the full-envelope target and \(\theta_\gamma\) for its marginal maximizer, the exact bound is

\[\left(\mathbb E\|\theta_n-\theta_\gamma\|^2\right)^{1/2}\le\sqrt{\frac{d_\theta}{Nm_\gamma}}+(1-m_\gamma\gamma)^nW_2(\rho_0,\Pi_{N,\gamma})+\frac{33}{20}\frac{\gamma^{-1}}{m_\gamma}\sqrt{d_x+d_\theta/N}\sqrt\gamma.\]

Prove full-envelope convexity and differentiability, identify the exact algorithm law, apply the sharp IPLA theorem with gradient constant \(\gamma^{-1}\), and use the already proved admissible-step inequality. This gives every displayed term from the actual chain.




Here the whole joint potential is smoothed, so the relevant modulus is \(m_\eta=m/(1+m\eta)\), rather than \(m\). The score constant of the full envelope is \(1/\eta\). The exact proximal/noise update identity reduces PIPULA to IPLA for that envelope, after which the theorem applies the sharp concentration and Gaussian discretization estimates with these constants. The admissible-step proof checks the actual inequality \(\gamma\le2/(m_\eta+1/\eta)\); substituting the unsmoothed modulus would give a different and unjustified step restriction.




\textbf{Correspondence with the displayed declarations.}


\begin{itemize}

\item \texttt{pipula\_parameter\_rms\_smoothed}: A complete finite-time bound for PIPULA around the maximizer of the smoothed marginal. The original potential may be nonsmooth and need not be globally Lipschitz. All envelope properties, the time-step restriction, normalization, finite transport moments, and equilibrium concentration are proved from the displayed assumptions and the actual Gaussian innovations.
\end{itemize}


\noindent Lean declarations: \texttt{pipula\_parameter\_rms\_smoothed}.

\end{document}
